\subsection{绝对正则代数的刻画}

命题 8. 设 \(k\) 为域，A 为本质有限型 \(k\)-代数。以 I 表示 \(k\)-代数同态 \(\mu: A \otimes_k A \to A\)（满足 \(\mu(x \otimes y) = xy\)，对 A 中的 \(x\) 与 \(y\)）的核。下列条件等价：

\begin{enumerate}
\def\labelenumi{(\roman{enumi})}
\item
  k-代数 A 绝对正则；
\item
  对每个正则 k-代数 C，环 A ⊗k C 正则；
\item
  环 A ⊗k A 正则；
\item
  A ⊗k A 的理想 I 是完全割的 (§ 5, 第 1 节, 定义 1)。
\end{enumerate}

以 B 表示环 A \(\otimes_{k}\) A，并赋予它由同态 \(\rho: \mathrm{A} \to \mathrm{A} \otimes_{k} \mathrm{A}\)（其中 \(\rho(x) = x \otimes 1\)）导出的 A-代数结构；于是 \(\mu\) 是 A-代数同态，并且通过过渡到商，诱导出从 B/I 到 A 上的同构。

(i)⇒(ii)：由命题 7 得出。

(ii)⇒(iii)：只需对 C=k 再对 C=A 应用 (ii)。

(iii)⇒(iv)：A-模 B 是自由的，从而忠实平坦。若环 B 正则，则 A 正则 (§ 4, 第 5 节, 命题 8)；于是理想 I 是完全割的 (§ 5, 第 3 节, 命题 2)。

(iv)⇒(i)：设理想 I 是完全割的，先证明 A 正则。设 m 为 A 的极大理想，设 \(\nu:(\mathrm{A}/\mathfrak{m})\otimes_{k}\mathrm{A}\to\mathrm{A}/\mathfrak{m}\) 为由 \(\mu\) 导出的同态。极大理想 n=Ker \(\nu\) 等于 \(I((\mathrm{A}/\mathfrak{m})\otimes_{k}A)\)；把 § 5 第 6 节的命题 6 应用于 A-代数 \(A'=A/\mathfrak m\)，可见理想 n 在 \((A/\mathfrak{m})\otimes_{k}A\) 中是完全割的。于是 (§ 5, 第 3 节, 命题 3)，局部环 \(((\mathrm{A}/\mathfrak{m})\otimes_{k}A)_{n}\) 正则。以 \(j:\mathrm{A}\to(\mathrm{A}/\mathfrak{m})\otimes_{k}\mathrm{A}\) 表示同态 \(x\mapsto1\otimes x\)；由于 \(\nu\circ j\) 是从 A 到 A/m 的典范同态，有 \(j^{-1}(\mathfrak{n})=\mathfrak{m}\)。于是 j 延拓为从 \(A_{m}\) 到 \(((\mathrm{A}/\mathfrak{m})\otimes_{k}\mathrm{A})_{n}\) 的局部环的局部同态，这使后者成为忠实平坦的 \(A_{m}\)-模。由 § 4 第 5 节的命题 8，环 \(A_{m}\) 因此正则。我们至此证明了 A 正则。

现在设 \(k'\) 为 \(k\) 的扩张。映射 \(\mu': A_{(k')} \otimes_{k'} A_{(k')} \to A_{(k')}\)（由 \(A_{(k')}\) 的乘法导出）的核正是 \(IA_{(k')}\)；因此它在 \(A_{(k')}\) 中是完全割的 (§ 5, 第 6 节, 命题 6)。于是 \(k'\)-代数 \(A_{(k')}\) 满足条件 (iv)；由刚才所见，它正则，这就证明了 (i)。

回忆 (A, III, 第 133–134 页)：商 I/I² 赋予由 ρ 诱导的 A-模结构后，记作 Ωₖ(A)，称为 A 的 k-微分模。当 k-代数 A 本质有限型时，环 A ⊗ₖ A 是 Noether 的，于是 A-模 Ωₖ(A) 有限生成。

以 \(d_{A/k}\) 或简称 \(d\) 表示一个 \(k\)-线性映射（自 A 到 \(\Omega_k(A)\)），它把 A 的元素 \(x\) 对应为 \(x \otimes 1 - 1 \otimes x\) 在 \(\Omega_k(A)\) 中的类。映射 \(d\) 是 \(k\)-导子；对每个 A-模 M 与每个 \(k\)-导子 D: A → M，存在唯一的 A-线性映射 \(g: \Omega_k(A) \to M\) 使 D = g ∘ d (同上引文, 命题 18)。

若 S 是 A 的乘性子集，则典范 \(S^{-1}A\)-线性映射 (同上引文, 第 136 页)

\[
\mathrm{S} ^ {- 1} \Omega_ {k} (\mathrm{A}) \rightarrow \Omega_ {k} (\mathrm{S} ^ {- 1} \mathrm{A})
\]

是双射的：事实上，只需验证：对每个 S \(^{-1}\) A-模 M，每个 k-导子 D : A → M 都唯一扩张为 k-导子 D : S \(^{-1}\) A → M；这由同上引文第 123 页命题 5 的论证得出。

设 \(k\) 为域，A 为本质有限型局部 \(k\)-代数。令 \(n = \dim(A) + \deg.\mathrm{tr}_k(\kappa_A)\)。设 B 为有限生成 \(k\)-代数，\(\mathfrak{q}\) 为 B 的素理想，使 \(k\)-代数 A 同构于 \(\mathbf{B}_{\mathfrak{q}}\)（\(cf.\ n^{\circ} 1\)，例 2）；我们有 \(n = \dim_{\mathfrak{q}}(\mathbf{B})\) (VIII, § 2, 第 4 节, 定理 3 的推论 5)。由同上引文定理 3, c) 与推论 5，还有

\[
n = \sup \deg . \operatorname{tr} _ {k} (\kappa (\mathfrak {p}))
\]

其中 p 取遍 A 的极小素理想的有限族。若 A 是整的，则 n 就是 A 的分式域的超越次数。

定理 1. 设 \(k\) 为域，A 为本质有限型局部 \(k\)-代数。令 \(n = \dim(A) + \deg.\mathrm{tr}_{k}(\kappa_{A})\)。

\begin{enumerate}
\def\labelenumi{\alph{enumi})}
\item
  有 \([\kappa_{A} \otimes_{A} \Omega_k(A): \kappa_{A}] \geqslant n\)。
\item
  下列条件等价：
\end{enumerate}

\begin{enumerate}
\def\labelenumi{(\roman{enumi})}
\item
  k-代数 A 绝对正则；
\item
  A-模 \(\Omega_{k}(\mathrm{A})\) 是秩为 \(n\) 的自由模；
\end{enumerate}

\[
\left[ \kappa_ {A} \otimes_ {A} \Omega_ {k} (A): \kappa_ {A} \right] = n
\]

考察断言

(iii′) 有 \([\kappa_{\mathrm{A}} \otimes_{A} \Omega_k(\mathrm{A}):\kappa_{\mathrm{A}}] \leqslant n\)；

显然 (ii) 蕴含 (iii)，且 (iii) 蕴含 (iii′)。为证明定理，只需证明蕴含关系 (i) \(\Rightarrow\) (ii) 与 (iii′) \(\Rightarrow\) (i)。

选取有限生成 \(k\)-代数 B 及 B 的素理想 \(\mathfrak{q}\)，使 A 同构于 \(\mathrm{B}_{\mathfrak{q}}\) (\(n^{\circ} 1\), 例 2)。我们有 \(\dim_{\mathfrak{q}}(\mathrm{B}) = n\)；把 B 换成环 \(\mathrm{B}_f\)（其中 \(f \in \mathrm{B} - \mathfrak{q}\)），可以要求 B 的谱连通且维数为 \(n\)，以下即作此假设。

\begin{enumerate}
\def\labelenumi{(\roman{enumi})}
\tightlist
\item
  \(\Rightarrow\) (ii)：设 \(k\)-代数 A 绝对正则。如同命题 8 的证明，以 \(\mu : A \otimes_k A \to A\) 与 \(\nu : \kappa_A \otimes_k A \to \kappa_A\) 表示由 A 的乘法导出的同态；令 I = \(\operatorname{Ker}(\mu)\)，\(\mathfrak{n} = \operatorname{Ker}(\nu)\)。理想 I 是完全割的；理想 \(\mathfrak{n}\) 极大，且局部环 (\(\kappa_A \otimes_k A\) ) \(_{\mathfrak{n}}\) 正则 (同上引文)。由 § 5 第 2 节的定理 1，有限型 \(A\)-模 I/I²——它正是 \(\Omega_k(A)\)——是投射的，从而是自由的。但依 § 5 第 6 节的注，理想 \(\mathfrak{n}\) 与 \(\kappa_A \otimes_A 1\) 恒同，于是 \(\mathfrak{n}/\mathfrak{n}^2\) 同构于 \(\kappa_A \otimes_A \Omega_k(A)\)。因此有
\end{enumerate}

\[
\operatorname{rg} \left(\Omega_ {k} (\mathrm{A})\right) = \left[ \kappa_ {\mathrm{A}} \otimes_ {k} \Omega_ {k} (\mathrm{A}): \kappa_ {\mathrm{A}} \right] = \left[ \mathfrak {n} / \mathfrak {n} ^ {2}: \kappa_ {\mathrm{A}} \right] = \dim \left(\kappa_ {\mathrm{A}} \otimes_ {k} \mathrm{A}\right) _ {\mathfrak {n}}.
\]

以 B′ 表示环 \(\kappa_{A}\otimes_{k}B\)，以 S 表示 B-q 在 B′ 中的像。环 \(\kappa_{A}\otimes_{k}A\) 与 B′ \(\otimes_{B}B_{q}\) 恒同，亦即与 S \(^{-1}\) B′ 恒同。因此存在一个不与 S 相交的极大理想 \(q'\)（属于 \(B'\)），使得 \(n = S^{-1}q'\) 成立；于是环 \((\kappa_A \otimes_k A)_n\) 同构于 \(B_{q'}'\)，且有

\[
\dim (\kappa_ {A} \otimes_ {k} A) _ {\mathfrak {n}} = \dim (\kappa_ {A} \otimes_ {k} B) _ {\mathfrak {q} ^ {\prime}} = \dim_ {\mathfrak {q} ^ {\prime}} (\kappa_ {A} \otimes_ {k} B).
\]

由于 n 包含 \(\kappa_{\mathrm{A}} \otimes_{k} qB_{q}\)，\(q'\) 包含理想 \(\kappa_{\mathrm{A}} \otimes q\)（\(B'\) 中的），于是它在 B 中的逆像包含 \(q\)；该逆像等于 \(q\)，因为 \(q'\) 不与 S 相交。由 VIII, § 2, 第 4 节, 命题 5 的推论，有

\[
\dim_ {\mathfrak {q} ^ {\prime}} \left(\kappa_ {A} \otimes_ {k} B\right) = \dim_ {\mathfrak {q}} (B) = n,
\]

这就证明了 (ii)。

(iii′) \(\Rightarrow\) (i)：设 \([\kappa_{\mathrm{A}}\otimes_{\mathrm{A}}\Omega_{k}(A):\kappa_{\mathrm{A}}]\leqslant n\)，即 \([\kappa (\mathfrak{q})\otimes_{\mathrm{B}}\Omega_{k}(\mathrm{B}):\kappa (\mathfrak{q})]\leqslant n\)，我们来证明 \(k\)-代数 B 绝对正则。设 \((x_{1},\ldots ,x_{n})\) 为 B 的元素序列，使 \(1\otimes dx_1,\dots ,1\otimes dx_n\) 生成 \(\kappa (\mathfrak{q})\)-向量空间 \(\kappa (\mathfrak{q})\otimes_{\mathrm{B}}\Omega_{k}(\mathrm{B})\)。把 B 换成 \(\mathrm{B}_f\)（\(f\) 为 \(\mathrm{B} - \mathfrak{q}\) 中适当选取的元素），可设 \(dx_{1},\ldots ,dx_{n}\) 生成 B-模 \(\Omega_k(\mathrm{B})\)。由 Nakayama 引理与 II, § 5, 第 1 节, 命题 2，设 \(\bar{k}\) 为 \(k\) 的代数闭扩张。我们只需证明 \(\bar{k}\)-代数 \(\mathrm{B}_{(\bar{k})}\) 正则，因为这将蕴含 B 绝对正则 (第 4 节命题 7 的推论 4)。对 \(\mathrm{B}_{(k)}\) 的每个典范分支 C，微分 \(d(1\otimes x_i)\) 生成 C-模 \(\Omega_{\bar{k}}(\mathrm{C})\) (A, III, § 10, 第 12 节, 命题 20)。于是定理由下面两个引理得出：

引理 4.—设 \(k\) 为域，K 为 \(k\) 的代数扩张。设 B 为谱连通的 \(k\)-代数。对 \(\mathrm{B}_{(\mathrm{K})}\) 的每个典范分支 C，有 \(\dim(\mathrm{C}) = \dim(\mathrm{B})\)。

先处理扩张 K 对 \(k\) 有限的情形。B-模 C 是自由模的直和项，从而是有限生成投射模。由于函数 \(\mathfrak{p} \mapsto \mathrm{rg}_{\mathfrak{p}}\) 在 \(\operatorname{Spec}(B)\) 上取常值 (II, § 5, 第 3 节)，B-模 C 的支撑集等于 \(\operatorname{Spec}(B)\)，于是 B-模 C 忠实平坦。这蕴含 \(\dim(C) \geqslant \dim(B)\) (VIII, § 2, 第 1 节, 命题 2)，从而得到所欲的等式，因为有 \(\dim(C) \leqslant \dim(B_{(K)}) = \dim(B)\) (VIII, § 2, 第 4 节, 命题 5 的推论)。

对一般情形，设 \(e\) 为 \(B_{(K)}\) 的幂等元，使 \(C = B_{(K)}/B_{(K)}e\)。存在一个有限次中间扩张 \(K'\)（在 \(K\) 上），使 \(e\) 是某个幂等元 \(e'\)（它属于 \(B_{(K')}\)）的像。令 \(C' = B_{(K')} / B_{(K')}e'\)。环 \(C' \otimes_{K'} K\) 与 \(C\) 恒同，且 \(C'\) 是 \(B_{(K')}\) 的一个典范分支。由已处理的情形有 \(\dim(C') = \dim(B)\)，又由同上引文有 \(\dim(C') = \dim(C)\)，引理得证。

引理 5.—设 \(k\) 为代数闭域，A 为谱连通的有限型 \(k\)-代数，\(n\) 为整数，\((x_{1},\ldots,x_{n})\) 为 A 的元素的一个有限序列。假设 A 的维数为 \(n\)，且微分 \(dx_{1},\ldots,dx_{n}\) 生成 A-模 \(\Omega_{k}(\mathrm{A})\)。则环 A 是整的且正则的，且 \(A\)-模 \(\Omega_{k}(\mathrm{A})\) 是自由的，以 \((dx_{1},\ldots,dx_{n})\) 为基。

设 \(\mathfrak{m}\) 为 A 的极大理想且 \(\dim(A_{\mathfrak{m}})=n\)。有 \([A/\mathfrak{m}:k]=1\) (V, § 3, 第 3 节, 命题 1 (iii))，于是 \(A=\mathfrak{m}\oplus k1_{A}\)。设 \(p\) 与 \(q\) 为相应的射影

对 \(a\) 与 \(b\)（\(A\) 中的元素），有

\[
a b = \left(p (a) q (b) + q (a) p (b) + p (a) p (b), q (a) q (b)\right),
\]

从而 \(p(ab) \equiv p(a)q(b) + q(a)p(b) \pmod{\mathfrak{m}^{2}}\)。于是，映射 \(\delta: A \to \mathfrak{m}/\mathfrak{m}^{2}\)（它把 A 的每个元素 x 对应为 \(p(x)\) 模 \(\mathfrak{m}^{2}\) 的类）是 A 到 k-向量空间 \(\mathfrak{m}/\mathfrak{m}^{2}\) 的 k-导子。因此存在 A-线性映射 \(\phi: \Omega_{k}(A) \to \mathfrak{m}/\mathfrak{m}^{2}\)，使 \(\delta(x) = \phi(dx)\) 对每个 \(x \in A\) 都成立。由于 \(\delta\) 是满射，诸 \(\phi(dx_{i})\) 生成 A/m-向量空间 \(\mathfrak{m}/\mathfrak{m}^{2}\)，且有 \([\mathfrak{m}/\mathfrak{m}^{2}: A/\mathfrak{m}] \leqslant n = \dim(A_{\mathfrak{m}})\)。由此推出 \(A_{\mathfrak{m}}\) 正则，且诸 \(dx_{i}\) 的像构成 A/m-向量空间 A/m \(\otimes_{A} \Omega_{k}(A)\) 的一组基。

现在证明环 A 是整的且正则的。存在 A 的极小素理想 q 使 \(\dim(A/q) = n\)。对 A 的每个包含 q 的极大理想 m，有 \(\dim(A_{\mathfrak{m}}) = n\) (VIII, § 2, 第 4 节, 定理 3 的推论 2)，于是由刚才所见，\(A_{\mathfrak{m}}\) 正则。特别地，\(A_{\mathfrak{m}}\) 是整的，这蕴含 \(qA_{\mathfrak{m}} = 0\)。由于这对 V(q) 的所有极大理想 m 都成立，我们推出 \(\operatorname{Supp}(\mathfrak{q}) \cap V(\mathfrak{q}) = \varnothing\)。但 \(\operatorname{Spec}(A)\) 连通，V(q) 非空，且有 \(\operatorname{Supp}(\mathfrak{q}) \cup V(\mathfrak{q}) = \operatorname{Spec}(A)\) (II, § 4, 第 4 节, 命题 16)。由此推出 \(\operatorname{Supp}(\mathfrak{q}) = \varnothing\)，从而 q = 0，这意味着 A 是整环。于是对 A 的每个极大理想 m 都有 \(\dim(A_{\mathfrak{m}}) = n\)；应用证明的第一部分，推出 A 正则。

最后，设 \(\sum_{i=1}^{n} a_i dx_i = 0\) 为诸 \(dx_i\) 之间的一个系数在 \(A\) 中的线性关系。若诸 \(a_i\) 不全为零，则存在指标 \(i\) 与 A 的极大理想 \(\mathfrak{m}\)，使 \(a_i\) 不属于 \(\mathfrak{m}\) (V, § 3, 第 3 节, 命题 1, (iii) 与 (iv))；但这与上面已证的事实矛盾：诸 \(dx_i\) 在 (A/m) \(\otimes_{\mathrm{A}} \Omega_k(\mathrm{A})\) 中的类线性无关。

例.--- 当 A 是 k 的有限生成扩张时，结合第 4 节的例 2，定理 1 重新给出 A, V, 第 128 页的推论 1。

推论 1. - 设 \(k\) 为域，\(A\) 为本质有限型 \(k\)-代数。素理想 \(\mathfrak{p}\)（属于 \(\operatorname{Spec}(A)\)）中，使 \(k\)-代数 \(\mathrm{A}_{\mathfrak{p}}\) 绝对正则的那些构成 \(\operatorname{Spec}(A)\) 的一个开子集。

可以假设 \(k\)-代数 A 是有限生成的。于是所考虑的集合由这样的素理想 \(\mathfrak{p}\) 组成：\([\kappa(\mathfrak{p}) \otimes_k \Omega_k(A) : \kappa(\mathfrak{p})] \leqslant \dim_{\mathfrak{p}}(A)\)。而函数 \(\mathfrak{p} \mapsto \dim_{\mathfrak{p}}(A)\) 依定义是下半连续的，函数 \(\mathfrak{p} \mapsto [\kappa(\mathfrak{p}) \otimes_k \Omega_k(A) : \kappa(\mathfrak{p})]\) 是上半连续的 (由 Nakayama 引理与 II, § 5, 第 1 节, 命题 2)。

我们稍后将看到 (§ 7, 第 9 节, 定理 3 的推论 4)：在推论 1 的假设下，A 中使环 \(A_{p}\) 正则的素理想 p 组成的集合在 \(\operatorname{Spec}(A)\) 中是开的。

推论 2.--- 设 \(k\) 为域，A 为本质有限生成的 \(k\)-代数。为使 A 绝对正则，必须且只需 A-模 \(\Omega_{k}(A)\) 是投射的，且对 A 的每个极小素理想 \(\mathfrak{q}\)，\(k\)-代数 \(\mathrm{A}_{\mathfrak{q}}\) 都是可分的。

设 A 绝对正则。对 A 的每个素理想 p，k-代数 \(A_{p}\) 绝对正则，于是 \(A_{p}\)-模 \(\Omega_{k}(A_{p})\) 是自由的 (定理 1)；由此推出 A-模 \(\Omega_{k}(A)\) 是投射的 (II, § 5, 第 2 节, 定理 1)。此外，对 A 的每个极小素理想 q，局部 k-代数 \(A_{q}\) 是 Artin 的且绝对正则，从而是一个域，即 k 的可分扩张 (第 4 节, 例 2)。

反之，设 A-模 \(\Omega_{k}(\mathbf{A})\) 是投射的，且 \(k\)-代数 \(\mathbf{A}_{\mathfrak{q}}\) 对每个极小素理想 \(\mathfrak{q}\)（\(\mathbf{A}\) 的）都可分。设 \(\mathfrak{p}\) 为 \(\mathbf{A}\) 的素理想，\(\mathfrak{q}\) 为 \(\mathbf{A}\) 的一个极小素理想（含于 \(\mathfrak{p}\)）。由于 \(\mathbf{A}_{\mathfrak{p}}\)-模 \(\Omega_{k}(\mathbf{A})_{\mathfrak{p}}\) 是自由的 (II, § 3, 第 2 节, 命题 5 的推论 2)，有

\[
[ \kappa (\mathfrak {p}) \otimes_ {\mathrm{A}} \Omega_ {k} (\mathrm{A}): \kappa (\mathfrak {p}) ] = [ \kappa (\mathfrak {q}) \otimes_ {\mathrm{A}} \Omega_ {k} (\mathrm{A}): \kappa (\mathfrak {q}) ].
\]

\(k\)-代数 \(\mathrm{A}_{\mathfrak{q}}\) 是 Artin 的且可分的，从而绝对正则 (\(n^{\circ}\) 4, 例 2)。定理 1 蕴含

\[
[ \kappa (\mathfrak {q}) \otimes_ {\mathrm{A}} \Omega_ {k} (\mathrm{A}): \kappa (\mathfrak {q}) ] = \deg . \operatorname{tr} _ {k} (\kappa (\mathfrak {q})).
\]

于是由定理 1 推出 \(k\)-代数 \(\mathrm{A}_{\mathfrak{p}}\) 绝对正则，推论得证。

注.--- 设 k-代数 A 绝对正则。于是 A 的全分式环 F 与诸 \(A_{q}\) 的积恒同，其中 q 取遍 A 的极小素理想组成的集合 (IV, § 2, 第 5 节, 命题 10)；因此它是可分 k-代数。反之，设 F 是可分 k-代数；对 A 的每个极小素理想 q，k-代数 \(A_{q}\) 是 F 的一个分式环 (IV, § 1, 第 1 节, 命题 2 的推论 3 与第 3 节, 命题 7 的推论 1)，从而是可分 k-代数 (第 4 节, 命题 6)。此外，若 A-模 \(\Omega_{k}(A)\) 是投射的，则由推论 2 推出 k-代数 A 绝对正则。

推论 3.--- 设 \(k\) 为特征 0 的域，A 为本质有限型 \(k\)-代数。为使 A 正则，必须且只需 A-模 \(\Omega_{k}(\mathrm{A})\) 是投射的。

由推论 2，只需证明：若 Artin 局部 \(k\)-代数 A 使 \(\Omega_k(A)\) 为自由 A-模，则 A 是域。而由 IX, § 3, 第 3 节, 定理 1，存在 A 的子域 K 使 \(A = K \oplus \mathfrak{m}_A\)。设 \(\delta: A \to \mathfrak{m}_A / \mathfrak{m}_A^2\) 为 A 到 \(\mathfrak{m}_A\) 上的投影与典范映射 \(\mathfrak{m}_A \to \mathfrak{m}_A / \mathfrak{m}_A^2\) 的复合。如引理 4 的证明那样，可验证 \(\delta\) 是 k-导子。但 A 到某个 A-模的每个 \(k\)-导子都零化 \(\mathfrak{m}_A\)：事实上，设 \(x \in \mathfrak{m}_A\)，选取整数 \(n \geqslant 1\) 使 \(x^{n-1} \neq 0\) 且 \(x^n = 0\)。这在 A-模 \(\Omega_k(A)\) 中蕴含关系 \(nx^{n-1} dx = 0\)，从而 \(dx = 0\)（因为 \(nx^{n-1} \neq 0\)）。于是 \(\delta(\mathfrak{m}_A) = 0\)，从而 \(\mathfrak{m}_A = \mathfrak{m}_A^2\)，最终得 \(\mathfrak{m}_A = \{0\}\)。

